% European Heart Journal: neutral adaptation of an inspected advance article.
\PassOptionsToPackage{no-math}{fontspec}
\documentclass[UTF8,fontset=none,scheme=plain]{ctexart}
\usepackage[paperwidth=595.276bp,paperheight=793.701bp,left=42.52bp,right=36.851bp,top=57bp,bottom=45bp,headheight=20bp,headsep=14bp,footskip=24bp]{geometry}
\newcommand{\JournalName}{European Heart Journal}
% Shared dependencies and neutral material. Compile with XeLaTeX.
% No publisher logos, paper text, data, figures, or download metadata are included.
\usepackage{amsmath,amssymb,graphicx,xcolor,tikz,fancyhdr,needspace,xurl,hyperref,balance,etoolbox}
% The legacy balance package omits the first split box's final baseline depth.
% Reserve that measured depth rather than hiding a CJK overfull-box warning.
\makeatletter
\patchcmd{\@BAlancecol}{\setbox\@leftcolumn=\vbox to \@tempdima}{\advance\@tempdima by \dp1\relax\setbox\@leftcolumn=\vbox to \@tempdima}{}{\PackageWarning{cns-template}{Balance depth patch unavailable; inspect final columns}}
\makeatother
\usetikzlibrary{arrows.meta,positioning}
\usepackage[protrusion=true,expansion=false]{microtype}
\IfFontExistsTF{Noto Serif CJK SC}{%
  \setCJKmainfont{Noto Serif CJK SC}%
}{%
  \setCJKmainfont[BoldFont=FandolSong-Bold.otf,ItalicFont=FandolKai-Regular.otf]{FandolSong-Regular.otf}%
}
\IfFontExistsTF{Noto Sans CJK SC}{%
  \setCJKsansfont{Noto Sans CJK SC}%
}{%
  \setCJKsansfont[BoldFont=FandolHei-Bold.otf]{FandolHei-Regular.otf}%
}
\setmainfont{TeX Gyre Termes}\setsansfont{TeX Gyre Heros}
\setmonofont{TeX Gyre Cursor}
\definecolor{JournalBlue}{HTML}{0097C9}
\definecolor{JournalRed}{HTML}{B83443}
\definecolor{ScienceRed}{HTML}{CD1021}
\definecolor{CiteBlue}{HTML}{0000EE}
\hypersetup{colorlinks=true,linkcolor=JournalBlue,urlcolor=JournalBlue,
  pdftitle={CNS论文汉化：中性版式示例},pdfauthor={示例模板；非期刊官方模板},
  pdfsubject={Neutral source-inspired Chinese layout demonstration; no real research claims}}
\urlstyle{same}\raggedbottom
\setlength{\parindent}{1em}\setlength{\parskip}{0pt}
\setlength{\emergencystretch}{2em}\tolerance=2200
\XeTeXlinebreaklocale "zh"\XeTeXlinebreakskip=0pt plus .1pt
\pagestyle{fancy}\fancyhf{}
\renewcommand{\headrulewidth}{0pt}\renewcommand{\footrulewidth}{0pt}
\newcommand{\PaperTitle}{中文题名占位：从组织结构到证据解释}
\newcommand{\SampleAuthors}{Author A\textsuperscript{1}, Author B\textsuperscript{2,*}}
\newcommand{\SampleAffiliations}{\textsuperscript{1}作者单位一（示例占位）。\quad
  \textsuperscript{2}作者单位二（示例占位）。\par
  *通信作者信息待填。DOI、收稿和出版日期均应从目标原文核实。}
\newcommand{\SampleAbstract}{本段是中性摘要占位文字，不陈述任何真实研究发现。它演示中文题名、作者信息、摘要和双栏正文之间的层级关系。实际使用时，应完整替换为目标论文的忠实中文译文，并保留研究对象、方法、结果和结论的限定条件。图注、参考文献、作者声明和来源信息均属于独立的覆盖项。模板中的图形仅说明排版结构，不表示生物学机制、实验结果或统计关系。}
\newcommand{\DemoCaption}{\textbf{图1. 中性证据链示意（占位图，无实验数据）。}
  \textbf{A}，输入节点表示需要核对的原文材料。
  \textbf{B}，处理节点表示逐段转换及复核。
  \textbf{C}，输出节点表示正文、图注和书目等独立交付部分。
  箭头仅表示本示例的阅读顺序，不代表因果关系、效应方向或数量变化。
  该图由模板自行绘制，所有节点均为文字占位；无样本量、误差线、统计检验或比例尺。
  实际论文图注应逐项保留原稿报告的样本量、单位、检验方法、限定及图形署名。}
\newcommand{\DemoVector}{%
 \resizebox{\linewidth}{!}{%
 \begin{tikzpicture}[font=\sffamily\small,node distance=18mm,
   box/.style={draw=black!70,rounded corners=2pt,fill=black!3,minimum width=26mm,minimum height=16mm,align=center},
   flow/.style={-{Stealth[length=2.5mm]},line width=.7pt}]
  \node[box] (a) {\textbf{A}\quad 输入\\原文材料};
  \node[box,right=of a] (b) {\textbf{B}\quad 处理\\转换与复核};
  \node[box,right=of b] (c) {\textbf{C}\quad 输出\\中文译稿};
  \draw[flow] (a)--(b);\draw[flow] (b)--(c);
  \node[below=11mm of a,align=center] {标题与摘要\\正文与图表};
  \node[below=11mm of b,align=center] {保留限定\\检查单位};
  \node[below=11mm of c,align=center] {可追溯引文\\明确原文信息};
  \node[below=29mm of b,font=\sffamily\footnotesize,text=black!65] {示意占位；不包含真实实验数据};
 \end{tikzpicture}}%
}
\newcommand{\SampleBibliography}{%
 \section*{参考文献（示例占位）}
 {\fontsize{7.6bp}{10bp}\selectfont\setlength{\parindent}{0pt}\interlinepenalty=10000
 \hypertarget{ref:1}{}1. Author A. (20XX). Example reference title. \textit{Journal placeholder}. 卷、页码与DOI待填。此条仅供演示引文跳转。\par\smallskip
 \hypertarget{ref:2}{}2. Author B. (20YY). Another example reference title. \textit{Journal placeholder}. 出版信息待填；不要把占位条目带入正式译稿。\par}
}
\newcommand{\SampleSourceInfo}{%
 \section*{原文信息}
 以下字段全部为示例占位：作者；英文原标题；期刊；年份；卷期；页码或文章号；DOI或来源链接。实际使用时必须逐项替换并核对。\par
 本文件是\JournalName 风格的非官方中文译稿样板，依据一次已核验文章版式进行结构化改编，不代表该期刊所有年代或文体的统一标准，也不保证像素级复刻。无真实论文内容或期刊标识图像。\par
}
% Deliberately neutral demo prose, not translated from a published paper.
\newcommand{\SampleOpening}{%
 \section*{引言：示例正文与引文}
 这是一段用于检查双栏阅读节奏的中性示例文字。段落中的术语、缩写和标点需要在实际项目中与源文保持一致。中文译稿应呈现原论文的完整论述，而不是把多段证据压缩成一段摘要。这里的引文只链接到本文件末尾明确标为占位的参考文献\SampleCiteOne。\par
 若原句仅表示相关、可能或提示，译文也应保留相同的证据强度。数字、单位、正负号、上下标和范围端点应逐项核对。原文存在内部矛盾时，应分别保留相应位置的陈述，再用简短译注说明；不能根据常识把不同的数字静默统一。\par
 \subsection*{内容层级与跨页衔接}
 此小标题演示二级层级。实际论文可能采用不同的小标题密度，模板不会替代对原文结构的识别。双栏文章需按左栏、右栏和跨页续句的顺序读取。若段落在页末中断，应在完成整句核对后接入下一页，避免把图中文字误当正文，或重复翻译同一段。\par
 图表具有独立的信息范围。正文引用到某幅图，并不意味着其图注可以省略。图注中的每个分面、比较对象、时间点、样本量及统计说明都应对应源文。英文图内标签可以保留，并在图旁或文末提供中文释义；实际采用哪种安排，需要依据图形尺寸和译文长度判断\SampleCiteTwo。\par
 \subsection*{来源与可维护数据}
 本段继续演示正文密度。可编辑文本和图注数据应与最终版面分开维护，以便修订时准确定位变化。若只调整版式，应能证明正文、图注和书目内容没有被意外改写。更换字体或纸张后，须重新检查断行、图注与图片的相对位置以及页脚可读性。\par
 文内引用与书目之间应有明确对应关系。这里的两条演示书目没有真实DOI，也没有虚构的实验信息。使用模板时可增删条目，但必须同时更新编号、链接目标和引用位置。页码只表示中文译稿自己的分页，不用于冒充原出版页码。\par
 \subsection*{结果叙述的演示占位}
 本节的文字仅用于呈现长段落在双栏页面中的流动，不报告真实结果。实际翻译应保留观察对象、实验条件和结局指标之间的关系。例如，原文分别描述体外实验与动物实验时，不能把两类结果合并为同一层面的结论；研究使用的对照、时间顺序和分析范围也必须保留。图表中的比较范围不应扩大为对整个人群或全部疾病阶段的判断。\par
 论述中的转折与限制同样属于正文。若研究先提出一种解释，再讨论另一种可能，译文应沿用这一推进顺序。作者明确写出的不确定性应能在中文中被识别；如果尚未确定某个现象的方向或机制，就不能改成已经建立的因果关系。对于术语近似但研究层级不同的概念，可在首次出现时保留缩写，后文保持统一。\par
 \subsection*{方法信息与可复核性}
 方法部分往往包含密集数字和专有名词。翻译时须区分浓度、总量、频率和处理时长，不能因为中文句式调整而改变参数的归属。试剂名称、软件版本和数据访问标识通常保留原样；上下限、排除规则和显著性标准则需要逐字符核对。这里没有提供任何真实实验参数，避免把模板占位误认为可以直接执行的方法。\par
 表格与补充材料的引用也需要保留原层级。正文对补充图的指向并不意味着模板中已经附有该补充图；若源文件只包含主文，应如实标注所交付的范围。图表编号、附录顺序和参考文献数量都应由目标文章决定，本示例中的结构不构成固定配额，也不要求所有论文采用同样的章节。\par
}
\newcommand{\SampleClosing}{%
 \section*{讨论：版式适配与范围}
 上方矢量图是模板原创的排版占位图，不来自任何论文。图形与完整演示图注被放在同一页，便于读者直接核对分面说明。真实项目应使用经过授权或用户提供的原图素材，并保留图注署名与许可信息。尺寸不足时，可适度缩放或调整栏宽，不以删除图注解决空间问题。\par
 模板仅记录一次已核验文章的可观察设计特征。即使文章出自同一期刊，不同年代、研究论文与综述、校正清样与正式版本也可能采用不同的页眉和引文形式。应以用户提供的实际PDF为准，不能仅凭期刊名称推断配色、字号或引用方式\SampleCiteOne。\par
 \subsection*{演示检查项}
 本样板用于展示可编译性、中文字形、标题层级、跨栏图注、页眉页脚和内部链接。图中没有实际数据，因此不应被当作科学插图引用。替换内容后，还需检查所有页的渲染结果，确认无截字、碰撞、异常空白和遗漏段落。\par
 \SampleBibliography
 \SampleSourceInfo
}


\usepackage{multicol,tabularx,array}
\definecolor{EHJRed}{HTML}{D6143D}
\definecolor{EHJBeige}{rgb}{.945,.945,.89}
\definecolor{EHJGrey}{rgb}{.9,.905,.905}
\setlength{\columnsep}{15.5bp}\raggedcolumns
\setlength{\parindent}{1em}\setlength{\parskip}{2bp}
\renewcommand{\normalsize}{\fontsize{9bp}{12.1bp}\selectfont}
\ctexset{section={format=\sffamily\bfseries\color{EHJRed}\fontsize{15bp}{18bp}\selectfont,beforeskip=11bp,afterskip=5bp},
 subsection={format=\sffamily\bfseries\fontsize{12.5bp}{15bp}\selectfont,beforeskip=9bp,afterskip=4bp}}
\hypersetup{colorlinks=true,linkcolor=black,urlcolor=blue,
 pdftitle={European Heart Journal风格中性中文样板},pdfauthor={非官方中性样板},
 pdfsubject={Original neutral layout example; no source article text or figures}}
\fancyhead[L]{\sffamily\bfseries\fontsize{9bp}{11bp}\selectfont\thepage}
\fancyhead[R]{\sffamily\fontsize{8bp}{10bp}\selectfont 示例作者等｜中性排版演示}
\renewcommand{\headrulewidth}{.5bp}
\renewcommand{\headrule}{\hbox to\headwidth{\color{EHJRed}\leaders\hrule height\headrulewidth\hfill}}
\fancyfoot[L]{\sffamily\fontsize{6.8bp}{9bp}\selectfont European Heart Journal｜原刊年份、卷期与页码按稿核验}
\fancyfoot[R]{\sffamily\fontsize{6.8bp}{9bp}\selectfont 非官方中性样板｜中文译页\thepage}
\newcommand{\EHJCite}[1]{\textsuperscript{\hyperlink{ref:#1}{#1}}}
\newcommand{\EHJDotRule}{\par\vspace{3bp}\noindent\leaders\hbox{\kern1bp.\kern1bp}\hfill\kern0pt\par\vspace{4bp}}
\newsavebox{\EHJPanelBox}
\newenvironment{ehjpanel}{\par\noindent\begingroup\setlength{\fboxsep}{8bp}\begin{lrbox}{\EHJPanelBox}\begin{minipage}{\dimexpr\linewidth-16bp\relax}\sffamily}{\end{minipage}\end{lrbox}\colorbox{EHJBeige}{\usebox{\EHJPanelBox}}\endgroup\par}
\newcommand{\EHJRow}[2]{\noindent\begin{tabularx}{\linewidth}{@{}>{\bfseries\raggedright\arraybackslash}p{.15\linewidth}X@{}}#1 & #2\end{tabularx}\par}
\newcommand{\EHJBody}{%
本段为原创中性排版文字，不陈述任何真实研究结论。正式全文译稿应把每个原文段落、方法、图注、表格、框文和声明分别列入内容清单，逐项保留比较对象、否定、数量、单位及证据范围。数字上标引文应连接到对应的英文条目，而非仅作为视觉装饰。\EHJCite{1}

图像和完整图注应接近其首次实质讨论的位置，沿原稿顺序与正文穿插。通栏图较小时，可以在图下或图上保留双栏文字；长图与长图注应重新适配，不能把全部主图集中到文章末尾，也不能为了维持英文页数而删减中文说明。

双栏阅读应先恢复跨栏、跨页续句，再组织中文段落。标题颜色、字体层级、浅色背景和点线只能辅助呈现，不能替代内容复核。中文长度改变分页后，仍须检查图表编号、书目目标和页眉页脚，确认没有遮挡、缺字或遗漏。}
\begin{document}\sffamily\normalsize
\fancypagestyle{ehjfirst}{\fancyhead{}\renewcommand{\headrulewidth}{0pt}}
\thispagestyle{ehjfirst}
\noindent\begin{minipage}[b]{.57\textwidth}
{\fontsize{14bp}{17bp}\selectfont European Heart Journal\par}
{\fontsize{8bp}{10bp}\selectfont 原刊年份、卷期、页码和DOI按稿核验\par}
\end{minipage}\hfill\begin{minipage}[b]{.40\textwidth}\raggedleft
{\bfseries\color{EHJRed}\fontsize{13bp}{16bp}\selectfont 转化研究（中性示例）\par}
{\fontsize{10bp}{12bp}\selectfont 学科信息占位\par}
\end{minipage}\par\vspace{9bp}
{\color{EHJRed}\hrule height .5bp}\vspace{20bp}
\noindent{\bfseries\fontsize{23bp}{28bp}\selectfont 中文题名占位：从内容清单到可追溯的完整译稿\par}\vspace{13bp}
\noindent{\bfseries\fontsize{13bp}{16bp}\selectfont 示例作者甲\textsuperscript{1,\ensuremath{\dagger}}；示例作者乙\textsuperscript{2,\ensuremath{\dagger}}；示例作者丙\textsuperscript{1,*}\par}\vspace{9bp}
{\fontsize{7.5bp}{10bp}\selectfont\noindent\textsuperscript{1}作者单位一（示例占位）；\textsuperscript{2}作者单位二（示例占位）。通信信息按原稿填写，不继承先前文章的地址或联系字段。\par}
\vspace{7bp}{\fontsize{8bp}{10bp}\selectfont 收稿、修回和接受日期：按所给版本逐项核验。\par}\vspace{16bp}
\begin{ehjpanel}
\textbf{\color{EHJRed}摘要（中性示例）}\par\vspace{8bp}
\EHJRow{背景与目的}{本段只演示结构化摘要的背景与目的，不报告真实研究。正式译稿应完整保留研究问题、对象、病理或理论语境和知识空白。不能根据题名补出原稿没有的动机，也不能用一段概括取代全文。正文、图注和补充材料实际是否提供，应分别记录；仅有链接不代表已经取得或完成翻译。}
\EHJDotRule
\EHJRow{方法}{本段演示方法摘要的标签窄栏与正文宽栏。实际稿应保留纳入范围、时间、样本量、实验模型、干预、测量、分析和统计方法。单位、上下标、缩写及基因符号均需核对；若源文存在实质矛盾，应区分忠实转述与译者说明。此处没有真实样本、数据或实验操作。}
\EHJDotRule
\EHJRow{结果}{本段演示结果摘要。中文表述应保留原文的比较方向、相关与因果强度、否定及不确定性。数值相同并不证明语义一致，还需核对分母、组别、观察时点和适用范围。图注中的统计定义和面板标记须完整呈现；本样板的图形及文字只是原创占位，不构成研究发现。}
\EHJDotRule
\EHJRow{结论}{本段演示结论摘要。真实结论应保持原稿的条件和谨慎程度，不把潜在用途写成已经验证的临床效益。排版以清楚可读为先；中文长度改变时可以自然续页，不能删去结论尾句。首页与后续结构化图形摘要可能内容相近，仍应作为不同内容单元完整保留。}
\end{ehjpanel}
\vfill\noindent\rule{\textwidth}{.4bp}\par
{\fontsize{7bp}{9bp}\selectfont *通信作者信息按原稿填写。\textsuperscript{\ensuremath{\dagger}}共同贡献关系按原稿核验。\par
非官方中文译稿样板，仅供学术交流。出版物版权归原作者及出版商所有。普通刊名替代品牌标识；无真实论文内容或原图。\par}

\clearpage
\begin{ehjpanel}
\textbf{\color{EHJRed}结构化图形摘要（中性示例）}\par\vspace{14bp}
\textbf{关键问题}\par
本段是问题陈述占位。正式稿应保留原文提出的核心问题、已有证据的局限和目标范围，不以译者自己的研究问题替换。\EHJDotRule
\textbf{关键发现}\par
本段是发现陈述占位。真实研究的实验对象、分析方法和比较方向均应逐项核对；不能将关联写成因果，也不能将图形摘要当成全部正文的替代品。\EHJDotRule
\textbf{核心启示}\par
本段是转化或解释意义占位。实际稿须保持潜在用途、适用范围与验证需求，不额外提供原文没有的建议。\par\vspace{12bp}
\begin{center}\colorbox{white}{\begin{minipage}{.91\linewidth}\vspace{7bp}\DemoVector\vspace{7bp}\end{minipage}}\end{center}
{\fontsize{8bp}{10.8bp}\selectfont 图形摘要为模板自绘的中性示意。输入、核对与输出只代表内容处理次序，不表示生物机制或实验结果；真实图形摘要的全部图内标签、箭头和完整图注均需核对。}\par
\EHJDotRule
\textbf{\color{EHJRed}关键词}\quad 中性示例；内容覆盖；图文对应；可追溯引文
\end{ehjpanel}\vspace{15bp}
\begingroup\setlength{\fboxsep}{7bp}\noindent\fcolorbox{black}{EHJGrey}{\begin{minipage}{\dimexpr\linewidth-15bp\relax}
{\bfseries\color{EHJRed}\fontsize{12.5bp}{15bp}\selectfont 转化展望（中性示例）\par}\smallskip
本框为独立内容单元的版面演示，不陈述任何真实治疗效果。实际稿须完整保留原文阐述的机制、干预对象、证据级别及临床转化所需的限定条件。正文以外的展望、编辑摘要和声明同样应列入覆盖清单，不能因其视觉上属于框文而遗漏。
\end{minipage}}\endgroup

\clearpage
\begin{multicols}{2}
\section*{正文与图文对应}
\EHJBody
\subsection*{层级与证据表达}
主节使用红色无衬线粗体，次节使用黑色无衬线粗体。正式稿应依原文层级组织，而不是为了视觉统一改写章节边界。\EHJCite{2}
\end{multicols}\vspace{5bp}
\noindent\begin{minipage}{\textwidth}
\begin{center}\begin{minipage}{.61\textwidth}\DemoVector\end{minipage}\end{center}
{\fontsize{8bp}{10.8bp}\selectfont\textbf{\color{EHJRed}图1．}\textbf{原创内容核对流程（中性示例，无实验数据）。}
\textbf{A}，输入节点表示标题、摘要、正文、图表与声明等需要核对的单元。\textbf{B}，处理节点表示恢复阅读顺序并逐项核对数据和范围。\textbf{C}，输出节点表示具有完整图注、可追溯书目和出处信息的译稿。箭头只演示处理次序，不代表因果效应、数量或生物机制。该图由模板自行绘制，未使用任何真实论文原图。}\par
\end{minipage}\vspace{8bp}
\begin{multicols}{2}
\subsection*{混排后的正文}
本段在通栏图注下恢复双栏，展示图与正文同页的阅读节奏。正式论文每幅主图都应靠近相应讨论，长图可独立安排，但不能把所有主图统一后附。原图中保留英文时，完整中文标签释义可单独附在文末。
\section*{参考文献（占位）}
{\fontsize{7.5bp}{10bp}\selectfont\setlength{\parindent}{0pt}\interlinepenalty=10000
\hypertarget{ref:1}{}1. Example author. Example source title. Journal and publication details to be verified. （中性占位，非真实文献。）\par\smallskip
\hypertarget{ref:2}{}2. Example author. Another example source title. DOI and source information to be verified. （中性占位，非真实文献。）\par}
\section*{声明（占位）}
致谢、资助、作者贡献、利益冲突、数据和许可说明均按原稿实际内容填写。本样板没有真实作者、研究数据或资助。
\section*{原文信息}
占位字段：作者、英文原题、期刊、年份、卷期、页码或文章号、DOI及明确列出的出版日期。Advance article与后续正式卷期可能不同；应以用户提供版本为依据，缺失字段不编造。

本文件为European Heart Journal风格的非官方中性样板。普通刊名、字体回退和中文分页为有意适配，不代表该刊所有文体，也不承诺逐像素一致。仅供学术交流，出版物版权归原作者及出版商所有。
\end{multicols}
\end{document}
